%% ma: L12-B08-0161
%% lop: 12
%% chuong: 2
%% bai: 8
%% dang: 12.8.51
%% loai: TN4
%% muc: TH
%% dap_an: "A"
%% nguon: "CĐ2.4.pdf giáo viên gửi 27/10/2026 (M07, Lê Minh Khôi – Nguyễn Anh Khôi, Chương II Chủ đề 4 Các bài toán về tâm tỉ cự), Câu 11"
%% kien_thuc: "Bài 8 (biểu thức tọa độ của phép toán vectơ, độ dài vectơ), Bài 6 (điểm thỏa hệ thức vectơ, quy tắc ba điểm)"
%% trang_thai: cho-duyet
%% ly_do: "Dạng mới 12.8.51 đề xuất, chờ giáo viên duyệt"
\begin{ex}
Trong không gian $Oxyz$, cho bốn điểm $A(1;4;0)$, $B(2;-1;-3)$, $C(-1;3;5)$ và $D(4;0;1)$. Gọi $M$ là điểm di động trong không gian sao cho biểu thức \[I=\left|\overrightarrow{MA}+\overrightarrow{MB}-2\overrightarrow{MC}-8\overrightarrow{MD}\right|\] đạt giá trị nhỏ nhất. Tọa độ của điểm $M$ là
\choice{\True $M\left(\dfrac{27}8;\dfrac38;\dfrac{21}8\right)$}{$M\left(\dfrac{27}8;\dfrac38;-\dfrac{21}8\right)$}{$M\left(-\dfrac{27}8;-\dfrac38;-\dfrac{21}8\right)$}{$M\left(\dfrac{27}8;-\dfrac38;-\dfrac{21}8\right)$}
\loigiai{Gọi $E$ thỏa $\overrightarrow{EA}+\overrightarrow{EB}-2\overrightarrow{EC}-8\overrightarrow{ED}=\vec 0\Rightarrow E=\dfrac{A+B-2C-8D}{-8}=\left(\dfrac{27}8;\dfrac38;\dfrac{21}8\right)$.

$I=8ME$ nhỏ nhất khi $M\equiv E$.}
\end{ex}
